\clearpage
\begingroup
\fontsize{9.5}{11.2}\selectfont
\raggedright
\setlength{\parskip}{0pt}
\begin{thebibliography}{99}
\setlength{\itemsep}{2pt}
\addcontentsline{toc}{section}{References}

\bibitem{proveit}
Vladimir Reshetnikov (repository owner).
\emph{ProveIt: Topology/UnknotRecognition and incoming research artifacts}.
Repository snapshot inspected on 9 October 2026, revision
\nolinkurl{8188525b70033dcfe7c51ea5ae2c8723ad0c0198}.
\href{https://github.com/VladimirReshetnikov/ProveIt/tree/8188525b70033dcfe7c51ea5ae2c8723ad0c0198/Topology/UnknotRecognition}{Pinned algorithm source tree};
\href{https://github.com/VladimirReshetnikov/ProveIt/tree/8188525b70033dcfe7c51ea5ae2c8723ad0c0198/docs/incoming}{pinned incoming directory}.

\bibitem{incoming-trace}
\emph{Single-Exponential Pachner Search and Exact Witness Reuse for Unknot
Recognition}.
Research report prepared for Vladimir Reshetnikov, 9 October 2026.
ProveIt incoming archive at the revision recorded above:
\href{https://github.com/VladimirReshetnikov/ProveIt/blob/8188525b70033dcfe7c51ea5ae2c8723ad0c0198/docs/incoming/unknot_trace_search_20261009.zip}{\nolinkurl{docs/incoming/unknot_trace_search_20261009.zip}}.
Relative to the extracted report root, article source:
\nolinkurl{article/article.tex};
connected-region enumeration and local commitment bounds are in
\nolinkurl{article/sections/regions.tex}.

\bibitem{incoming-active}
\emph{Active Normal Faces and Batched Geometric Descent:
Certified improvements for unknot recognition}.
Research continuation for Vladimir Reshetnikov, 9 October 2026.
ProveIt incoming archive at the revision recorded above:
\href{https://github.com/VladimirReshetnikov/ProveIt/blob/8188525b70033dcfe7c51ea5ae2c8723ad0c0198/docs/incoming/unknot_active_faces_batched_descent_20261009.zip}{\nolinkurl{docs/incoming/unknot_active_faces_batched_descent_20261009.zip}}.
Relative to the extracted report root, article source:
\nolinkurl{article/article.tex}.

\bibitem{incoming-envelope}
\emph{Minimum Envelopes and Complete Normal-Surface Sectors:
Linear ray bounds, exact enumeration, and independent coverage certificates}.
A research continuation for the ProveIt project, 9 October 2026.
ProveIt incoming archive at the revision recorded above:
\href{https://github.com/VladimirReshetnikov/ProveIt/blob/8188525b70033dcfe7c51ea5ae2c8723ad0c0198/docs/incoming/unknot_minimum_envelopes_20261009.zip}{\nolinkurl{docs/incoming/unknot_minimum_envelopes_20261009.zip}}.
Relative to the extracted report root, article source:
\nolinkurl{article/minimum_envelopes.tex}.

\bibitem{burton-he-counterexamples}
Benjamin A. Burton and Alexander He.
Finding large counterexamples by selectively exploring the Pachner graph.
In \emph{39th International Symposium on Computational Geometry (SoCG 2023)},
Leibniz International Proceedings in Informatics \textbf{258},
21:1--21:16, Schloss Dagstuhl--Leibniz-Zentrum f\"ur Informatik, 2023.
\href{https://doi.org/10.4230/LIPIcs.SoCG.2023.21}{doi:10.4230/LIPIcs.SoCG.2023.21}.
Expanded version consulted: \href{https://arxiv.org/abs/2303.06321v2}{arXiv:2303.06321v2},
7 March 2024, 37 pages; see Sections 2.2 and 3.3 for elementary moves
and object tracking.

\bibitem{kanj-xia}
Iyad Kanj and Ge Xia.
Flip distance is in FPT time $O(n+k\cdot c^k)$.
In \emph{32nd International Symposium on Theoretical Aspects of Computer
Science (STACS 2015)}, Leibniz International Proceedings in Informatics
\textbf{30}, 500--512, Schloss Dagstuhl--Leibniz-Zentrum f\"ur Informatik, 2015.
\href{https://doi.org/10.4230/LIPIcs.STACS.2015.500}{doi:10.4230/LIPIcs.STACS.2015.500}.
\href{https://drops.dagstuhl.de/entities/document/10.4230/LIPIcs.STACS.2015.500}{Publisher's open-access article}.

\bibitem{lackenby-talk}
Marc Lackenby.
\emph{Unknot recognition in quasi-polynomial time}.
Oxford topology seminar slides, March 2021, 34-page compressed handout.
\url{https://people.maths.ox.ac.uk/lackenby/quasipolynomial-talk-oxford-compressed.pdf}.
The main theorem stated on PDF page 7 has complexity
$2^{O((\log n)^3)}$; this reference is an author-level announcement and
proof outline, cited in that form.

\bibitem{lackenby-2026}
Marc Lackenby.
\emph{Incompressible surfaces, hierarchies and unknot recognition}.
Preprint, \href{https://arxiv.org/abs/2607.23350v1}{arXiv:2607.23350v1},
25 July 2026, 54 pages.
\href{https://arxiv.org/html/2607.23350v1}{Full text};
\href{https://people.maths.ox.ac.uk/lackenby/algorithm-incompressible-250726.pdf}{author's PDF}.
Section 9 and Proposition 9.1 give the iteration bound $L(g+1)^L$;
the section begins by explaining that some steps have not yet been specified
precisely enough to estimate their running time.  It is not cited here as a
paper establishing the quasi-polynomial bound announced in the 2021 slides.

\bibitem{nss}
Moni Naor, Leonard J. Schulman, and Aravind Srinivasan.
Splitters and near-optimal derandomization.
In \emph{Proceedings of the 36th Annual Symposium on Foundations of Computer
Science (FOCS 1995)}, 182--191, IEEE Computer Society Press, 1995.
\href{https://doi.org/10.1109/SFCS.1995.492475}{doi:10.1109/SFCS.1995.492475}.
\href{https://www.wisdom.weizmann.ac.il/~naor/PAPERS/splitters.pdf}{Authors' preliminary-version PDF},
Theorem 3(iii): the total-time perfect-hash-family construction used here.

\bibitem{color-coding}
Noga Alon, Raphael Yuster, and Uri Zwick.
Color-coding.
\emph{Journal of the ACM} \textbf{42}(4), 844--856, 1995.
\href{https://doi.org/10.1145/210332.210337}{doi:10.1145/210332.210337}.
\href{https://eccc.weizmann.ac.il/report/1994/009/}{Authors' ECCC report TR94-009},
12 December 1994.

\bibitem{gutin-color}
Gregory Gutin, Felix Reidl, Magnus Wahlstr\"om, and Meirav Zehavi.
Designing deterministic polynomial-space algorithms by color-coding
multivariate polynomials.
\emph{Journal of Computer and System Sciences} \textbf{95}, 69--85, 2018.
\href{https://doi.org/10.1016/j.jcss.2018.01.004}{doi:10.1016/j.jcss.2018.01.004}.
Version consulted: \href{https://arxiv.org/abs/1706.03698v2}{arXiv:1706.03698v2},
19 December 2017; Theorem 4 and Section 7 establish the deterministic
polynomial-delay splitter enumeration used for the bibliographic comparison.

\bibitem{regina}
Benjamin A. Burton, Ryan Budney, William Pettersson, et al.
\emph{Regina: Software for low-dimensional topology}.
1999--2025. Version 7.4 used for the independent triangulation controls.
\url{https://regina-normal.github.io/}.
\href{https://regina-normal.github.io/engine-docs/}{Engine API documentation}.
Software citation follows the project's recommended attribution; website
and documentation consulted on 9 October 2026.

\end{thebibliography}
\endgroup
